% GENERATED FILE. Edit canonical lessons, then run npm run generate:books.
% canonical-source-hash: fnv1a64:f2e80873023a9870
% canonical-lessons: KA-C311-anta-anisuttade, KA-C311-hagagi, KA-R311-opinion-and-result

\chapter{What I Think, and So}
\label{ch:ka-what-i-think-and-so}
\hlchaptermodality{311}

\begin{chapteropening}
\textbf{By the end of this chapter:} \emph{I can give my view with … anta nanage anisuttade and nanna abhiprāyadalli, ask for someone else's, and give a result with hāgāgi, adakkē and ādudarinda.}

Complete the last lesson of chapter 311: give a view with anta nanage anisuttade and its result with hāgāgi.
\end{chapteropening}

\section[\texorpdfstring{\kn{ಅಂತ} \kn{ನನಗೆ} \kn{ಅನಿಸುತ್ತದೆ} (anta nanage anisuttade)}{ಅಂತ ನನಗೆ ಅನಿಸುತ್ತದೆ (anta nanage anisuttade)}]{\kn{ಅಂತ} \kn{ನನಗೆ} \kn{ಅನಿಸುತ್ತದೆ} (anta nanage anisuttade) — what I think}

\label{lesson:KA-C311-anta-anisuttade}



\begin{quote}
Say \textquotedblleft{}to seem, to feel (that)\textquotedblright{} (\emph{anisu}) and the word that closes a quotation (\emph{anta}). Then \textquotedblleft{}clean\textquotedblright{} (\emph{svaccha}). Which letter sits small inside it? (\textbf{\kn{ಛ}}, \emph{cha}.)
\end{quote}

\subsection*{You'll want to know: … \kn{ಅಂತ} \kn{ನನಗೆ} \kn{ಅನಿಸುತ್ತದೆ}}
\begin{quote}
\textbf{\kn{ಈ} \kn{ಪುಸ್ತಕ} \kn{ಚೆನ್ನಾಗಿದೆ} \kn{ಅಂತ} \kn{ನನಗೆ} \kn{ಅನಿಸುತ್ತದೆ}.}
\end{quote}

\begin{quote}
\emph{ī pustaka cennāgide anta nanage anisuttade.} — \textquotedblleft{}I think this book is good.\textquotedblright{}
\end{quote}

The thought comes first, whole: \textbf{\kn{ಈ} \kn{ಪುಸ್ತಕ} \kn{ಚೆನ್ನಾಗಿದೆ}} (\emph{ī pustaka cennāgide}, \textquotedblleft{}this book is good\textquotedblright{}). \textbf{\kn{ಅಂತ}} closes it like a pair of quotation marks. Then \textbf{\kn{ನನಗೆ} \kn{ಅನಿಸುತ್ತದೆ}}: \textquotedblleft{}it seems to me\textquotedblright{}. Kannada does not say \textquotedblleft{}I think\textquotedblright{} here. The thought comes \textbf{to} you, so \textquotedblleft{}me\textquotedblright{} takes the \textbf{\kn{ನನಗೆ}} of \textbf{\kn{ನನಗೆ} \kn{ಕನ್ನಡ} \kn{ಗೊತ್ತು}}.

To ask for someone's view:

\begin{quote}
\textbf{\kn{ನಿಮಗೆ} \kn{ಏನು} \kn{ಅನಿಸುತ್ತದೆ}?}
\end{quote}

\begin{quote}
\emph{nimage ēnu anisuttade?} — \textquotedblleft{}What do you think?\textquotedblright{}
\end{quote}

For what you thought then, the \textquotedblleft{}it\textquotedblright{} past: \textbf{\kn{ಅನಿಸಿತು}}, \textquotedblleft{}it seemed\textquotedblright{}.

\begin{quote}
\textbf{\kn{ಅದು} \kn{ಕಷ್ಟ} \kn{ಅಂತ} \kn{ನನಗೆ} \kn{ಅನಿಸಿತು}.}
\end{quote}

\begin{quote}
\emph{adu kaṣṭa anta nanage anisitu.} — \textquotedblleft{}I felt that it was difficult.\textquotedblright{}
\end{quote}

\begin{grammarlens}[title={Grammar Lens: \kn{ನನ್ನ} \kn{ಅಭಿಪ್ರಾಯದಲ್ಲಿ}, in my opinion}]
\begin{quote}
\textbf{\kn{ನನ್ನ} \kn{ಅಭಿಪ್ರಾಯದಲ್ಲಿ} \kn{ಈ} \kn{ಹೋಟೆಲ್} \kn{ಒಳ್ಳೆಯದು}.}
\end{quote}

\begin{quote}
\emph{nanna abhiprāyadalli ī hōṭel oḷḷeyadu.} — \textquotedblleft{}In my opinion this hotel is good.\textquotedblright{}
\end{quote}

\textbf{\kn{ಅಭಿಪ್ರಾಯದಲ್ಲಿ}} is \textbf{\kn{ಅಭಿಪ್ರಾಯ}}, \textquotedblleft{}an opinion\textquotedblright{}, with \textbf{-\kn{ದಲ್ಲಿ}}, \textquotedblleft{}in\textquotedblright{}. \textbf{\kn{ಒಳ್ಳೆಯದು}} (\emph{oḷḷeyadu}, \textquotedblleft{}a good one; (it is) good\textquotedblright{}): \textbf{\kn{ಒಳ್ಳೆಯ}} goes in front of a noun, and at the end of a sentence it takes \textbf{-\kn{ದು}}.

In talk you will also hear \textbf{\kn{ನನ್ನ} \kn{ಪ್ರಕಾರ}} (\emph{nanna prakāra}, \textquotedblleft{}according to me\textquotedblright{}). \textbf{\kn{ಯೋಚಿಸು}} stays the verb for thinking something over; for a view given in passing, \textbf{\kn{ಅನಿಸುತ್ತದೆ}} is the everyday one.
\end{grammarlens}

\subsection*{Your turn}
Say these aloud:

\begin{itemize}
\raggedright
  \item \textquotedblleft{}I think this book is good\textquotedblright{} — \emph{ī pustaka cennāgide anta nanage anisuttade}
  \item \textquotedblleft{}What do you think?\textquotedblright{} — \emph{nimage ēnu anisuttade?}
  \item \textquotedblleft{}I felt that it was difficult\textquotedblright{} — \emph{adu kaṣṭa anta nanage anisitu}
  \item \textquotedblleft{}In my opinion this hotel is good\textquotedblright{} — \emph{nanna abhiprāyadalli ī hōṭel oḷḷeyadu}
\end{itemize}

\begin{tcolorbox}[breakable,colback=teal!4,colframe=teal!35!black,title={Before you move on}]
Where does the thought go? (\textbf{First, closed by }\kn{ಅಂತ}\textbf{.}) Who takes \textbf{\kn{ನನಗೆ}}? (\textbf{The one the thought comes to: me.}) \textquotedblleft{}What do you think?\textquotedblright{} (\textbf{\kn{ನಿಮಗೆ} \kn{ಏನು} \kn{ಅನಿಸುತ್ತದೆ}?}) \textquotedblleft{}In my opinion\textquotedblright{}? (\textbf{\kn{ನನ್ನ} \kn{ಅಭಿಪ್ರಾಯದಲ್ಲಿ}}.)
\end{tcolorbox}
\section[\texorpdfstring{\kn{ಹಾಗಾಗಿ} (hāgāgi)}{ಹಾಗಾಗಿ (hāgāgi)}]{\kn{ಹಾಗಾಗಿ} (hāgāgi) — so}

\label{lesson:KA-C311-hagagi}



\begin{quote}
Say \textquotedblleft{}I think this book is good\textquotedblright{}. (\emph{ī pustaka cennāgide anta nanage anisuttade}.) Then \textquotedblleft{}that is why\textquotedblright{} and \textquotedblleft{}therefore\textquotedblright{}. (\emph{adakkē}, \emph{ādudarinda}.)
\end{quote}

\subsection*{You'll want to know: \kn{ಹಾಗಾಗಿ}}
\begin{quote}
\textbf{\kn{ನಿನ್ನೆ} \kn{ಮಳೆ} \kn{ಬಂತು}. \kn{ಹಾಗಾಗಿ} \kn{ನಾನು} \kn{ಮನೆಯಲ್ಲೇ} \kn{ಇದ್ದೆ}.}
\end{quote}

\begin{quote}
\emph{ninne maḷe bantu. hāgāgi nānu maneyallē idde.} — \textquotedblleft{}Yesterday it rained. So I stayed at home.\textquotedblright{}
\end{quote}

The reason comes first, then \textbf{\kn{ಹಾಗಾಗಿ}}, then the result. Taken apart, it is \textbf{\kn{ಹಾಗೆ}}, \textquotedblleft{}that way\textquotedblright{}, and \textbf{\kn{ಆಗಿ}}, the \textbf{-\kn{ಇ}} form of \textbf{\kn{ಆಗು}}: \textquotedblleft{}that having happened\textquotedblright{}. \textbf{\kn{ಮನೆಯಲ್ಲೇ}} (\emph{maneyallē}, \textquotedblleft{}at home, nowhere else\textquotedblright{}) is \textbf{\kn{ಮನೆಯಲ್ಲಿ}}, \textquotedblleft{}at home\textquotedblright{}, with \textbf{-\kn{ಏ}}, which leans on the word: that very place.

\begin{quote}
\textbf{\kn{ನನಗೆ} \kn{ಜ್ವರ} \kn{ಇತ್ತು}. \kn{ಹಾಗಾಗಿ} \kn{ನಾನು} \kn{ಶಾಲೆಗೆ} \kn{ಹೋಗಲಿಲ್ಲ}.}
\end{quote}

\begin{quote}
\emph{nanage jvara ittu. hāgāgi nānu śālege hōgalilla.} — \textquotedblleft{}I had a fever. So I did not go to school.\textquotedblright{}
\end{quote}

\subsection*{You'll want to know: three words for 'so'}
\begin{itemize}
\raggedright
  \item \textbf{\kn{ಹಾಗಾಗಿ}} — the everyday \textquotedblleft{}so\textquotedblright{}, in talk and in writing
  \item \textbf{\kn{ಅದಕ್ಕೇ}} — \textquotedblleft{}that is why\textquotedblright{}, with weight on the reason: \textquotedblleft{}that is exactly why\textquotedblright{}
  \item \textbf{\kn{ಆದುದರಿಂದ}} — \textquotedblleft{}therefore\textquotedblright{}, in careful writing and in speeches
\end{itemize}

\begin{quote}
\textbf{\kn{ಬಸ್} \kn{ಬರಲಿಲ್ಲ}. \kn{ಅದಕ್ಕೇ} \kn{ನಾನು} \kn{ತಡವಾಗಿ} \kn{ಬಂದೆ}.}
\end{quote}

\begin{quote}
\emph{bas baralilla. adakkē nānu taḍavāgi bande.} — \textquotedblleft{}The bus didn't come. That's why I came late.\textquotedblright{}
\end{quote}

\textbf{\kn{ತಡವಾಗಿ}} (\emph{taḍavāgi}, \textquotedblleft{}late\textquotedblright{}). \textbf{\kn{ಆದುದರಿಂದ}} is built from \textbf{\kn{ಆಗು}} as well: \textquotedblleft{}from what happened\textquotedblright{}. All three stand at the start of the result, never in front of the reason.

\subsection*{Your turn}
Say these aloud:

\begin{itemize}
\raggedright
  \item \textquotedblleft{}Yesterday it rained. So I stayed at home.\textquotedblright{} — \emph{ninne maḷe bantu. hāgāgi nānu maneyallē idde.}
  \item \textquotedblleft{}I had a fever. So I did not go to school.\textquotedblright{} — \emph{nanage jvara ittu. hāgāgi nānu śālege hōgalilla.}
  \item \textquotedblleft{}The bus didn't come. That's why I came late.\textquotedblright{} — \emph{bas baralilla. adakkē nānu taḍavāgi bande.}
\end{itemize}

\begin{tcolorbox}[breakable,colback=teal!4,colframe=teal!35!black,title={Before you move on}]
\textquotedblleft{}So\textquotedblright{}, in talk? (\textbf{\kn{ಹಾಗಾಗಿ}}.) \textquotedblleft{}That is why\textquotedblright{}? (\textbf{\kn{ಅದಕ್ಕೇ}}.) \textquotedblleft{}Therefore\textquotedblright{}, in writing? (\textbf{\kn{ಆದುದರಿಂದ}}.) Which comes first, the reason or the result? (\textbf{The reason.}) What is \textbf{\kn{ಹಾಗಾಗಿ}} taken apart? (\textbf{\kn{ಹಾಗೆ}} + \textbf{\kn{ಆಗಿ}}: \textquotedblleft{}that having happened\textquotedblright{}.)
\end{tcolorbox}
\section[(review)]{Review — a view and its result}

\label{lesson:KA-R311-opinion-and-result}



\begin{quote}
Read \textbf{\kn{ಸ್ವಚ್ಛ}} aloud. (\emph{svaccha}.) Which letter sits small under \textbf{\kn{ಚ}}? (\textbf{\kn{ಛ}}.) Then \textquotedblleft{}What do you think?\textquotedblright{} (\emph{nimage ēnu anisuttade?})

\begin{itemize}
\raggedright
  \item \emph{Read it:} the story of the day the guests came, once more, aloud
\end{itemize}
\end{quote}

\subsection*{Your turn}
Say these aloud:

\begin{itemize}
\raggedright
  \item I think this hotel is good. In my opinion it is clean.
  \item I felt that the book was difficult.
  \item There was no bus. So I walked home.
  \item I had a fever yesterday. That's why I did not come.
  \item one thing you think about your town, and one result, with -\kn{ಅಂತ} \kn{ನನಗೆ} \kn{ಅನಿಸುತ್ತದೆ} and \kn{ಹಾಗಾಗಿ}
\end{itemize}

\begin{tcolorbox}[breakable,colback=teal!4,colframe=teal!35!black,title={Before you move on}]
\textquotedblleft{}I think that …\textquotedblright{}? (\textbf{… }\kn{ಅಂತ} \kn{ನನಗೆ} \kn{ಅನಿಸುತ್ತದೆ}\textbf{.}) \textquotedblleft{}In my opinion\textquotedblright{}? (\textbf{\kn{ನನ್ನ} \kn{ಅಭಿಪ್ರಾಯದಲ್ಲಿ}}.) \textquotedblleft{}So\textquotedblright{}? (\textbf{\kn{ಹಾಗಾಗಿ}}.) \textquotedblleft{}That is why\textquotedblright{}? (\textbf{\kn{ಅದಕ್ಕೇ}}.)
\end{tcolorbox}
